% BXH-switches.tex -- testing successive switches with BXH and BXLE.
% Build: xelatex BXH-switches.tex (twice).
\documentclass[11pt]{article}
\usepackage{fontspec}
\usepackage{polyglossia}
\setdefaultlanguage{english}
\setotherlanguage{hebrew}
\setmainfont{Times New Roman}
\newfontfamily\hebrewfont[Script=Hebrew]{Times New Roman}
\setmonofont{Consolas}[Scale=MatchLowercase]
\usepackage[margin=1in]{geometry}
\usepackage{amsmath}
\usepackage{booktabs}
\usepackage{tikz}
\usetikzlibrary{arrows.meta,positioning,calc}
\usepackage{fancyvrb}
\usepackage[hidelinks]{hyperref}

\title{Testing Successive Switches with BXH and BXLE}
\author{Seymour J. Metz (H: Shmuel \texthebrew{שמואל בן לייביש ולאה})}
\date{10 November 2026}

% One register cell for the diagrams: \cell{x}{y}{text}{fill}
\newcommand{\cell}[4]{%
  \node[draw, minimum width=7mm, minimum height=7mm, fill=#4,
        font=\ttfamily\small] at (#1,#2) {#3};}

\begin{document}
\maketitle

\begin{abstract}
A single BXH or BXLE instruction whose R1 and R3 fields name the same odd
general register both tests one bit of a string of switches and advances to
the next. A sequence of such instructions tests up to 30 switches in a
fullword with no shift, mask or test-under-mask instruction. The technique
is due to William Collier, who devised it in 1965 for System/360 and
described it in IBM Technical Report TR 00.1412-1~\cite{collier1966}. This
paper explains, from the definition of the instructions in the System/360
Principles of Operation~\cite{poo360}, why it works, and shows the layout of
the switch word that it requires.
\end{abstract}

\section{Origin}

William Collier devised the technique in 1965 at IBM Poughkeepsie, while
working on an operating system for System/360 that had to fit in 1K bytes;
by his account it saved some thirty bytes. He described it in IBM Technical
Report TR 00.1412-1, dated 22 June 1966~\cite{collier1966}, which is not
known to be available online. He recalled it in a note of 2 November 2022,
which the present author forwarded to the ASSEMBLER-LIST and IBM-MAIN
mailing lists in June 2024~\cite{metz2024,ibmmain2024}.

\section{The instructions}

BRANCH ON INDEX HIGH (BXH) and BRANCH ON INDEX LOW OR EQUAL (BXLE) are
RS-format instructions with operands $R_1$, $R_3$, $D_2(B_2)$. In the words
of the System/360 Principles of Operation~\cite[p.~66]{poo360}, ``an
increment is added to the first operand, and the sum is compared
algebraically with a comparand.'' The sum replaces the first operand whether
or not the branch is taken. BXH branches when the sum is high; BXLE branches
when it is low or equal~\cite[p.~66.1]{poo360}.

Three further details of the definition make the technique work:
\begin{enumerate}
\item The increment is in register $R_3$. The comparand register is odd, and
  is either $R_3 + 1$ or $R_3$ itself~\cite[p.~66.1]{poo360}. When $R_3$ is
  odd, therefore, one register supplies both the increment and the
  comparand.
\item ``Overflow caused by the addition is ignored and does not affect the
  comparison''; the comparison is of signed 32-bit values in two's
  complement~\cite[p.~66]{poo360}.
\item ``When the first operand and comparand locations coincide, the
  original register contents are used as the comparand''~\cite[p.~66]{poo360}.
\end{enumerate}

\section{The trick}

Let $R_1 = R_3 = r$, an odd register holding the value $x$. By the rules
above, the instruction computes the sum $x + x = 2x$, which is $x$ shifted
left one bit, with overflow ignored. It compares that sum with the
\emph{original} $x$, stores the sum in $r$, and branches (for BXH) when
$2x > x$. Figure~\ref{fig:one} shows one execution.

\begin{figure}[ht]
\centering
\begin{tikzpicture}[x=7mm,y=7mm]
  \node[anchor=east, font=\small] at (-0.8,0) {$x$ (before)};
  \foreach \b [count=\i from 0] in {s,b,?,?,?,?,?,?} {
    \ifnum\i=0 \cell{\i}{0}{\b}{gray!25}
    \else\ifnum\i=1 \cell{\i}{0}{\b}{yellow!40}
    \else \cell{\i}{0}{\b}{white}\fi\fi }
  \node[anchor=east, font=\small] at (-0.8,-2.4) {$2x$ (after)};
  \foreach \b [count=\i from 0] in {b,?,?,?,?,?,?,?} {
    \ifnum\i=0 \cell{\i}{-2.4}{\b}{yellow!40}
    \else \cell{\i}{-2.4}{\b}{white}\fi }
  \foreach \i in {1,...,7} {
    \draw[-{Stealth}, thick] (\i,-0.45) -- (\i-1,-1.95); }
  \node[font=\small, align=left, anchor=west] at (8.3,-1.2)
    {every bit moves one place left;\\
     $s$ (the old sign) is lost;\\
     $b$ becomes the new sign;\\
     the comparand is still the old $x$};
\end{tikzpicture}
\caption{One execution of \texttt{BXH r,r,target}: the register is doubled,
that is, shifted left one bit (shown at 8 bits; the machine uses 32).}
\label{fig:one}
\end{figure}

\section{Why it tests one bit}

Let $s$ be the sign bit of $x$ and $b$ the bit that follows it, the bit that
becomes the sign of $2x$. Table~\ref{tab:cases} takes the four cases.

\begin{table}[ht]
\centering
\begin{tabular}{ccll}
\toprule
$s$ & $b$ & $2x$ (overflow ignored) & is $2x > x$? \\
\midrule
0 & 0 & non-negative, no overflow & yes, provided $x > 0$ \\
0 & 1 & negative (overflow) & no: $2x < 0 \le x$ \\
1 & 0 & non-negative (overflow) & yes: $2x \ge 0 > x$ \\
1 & 1 & negative, no overflow & no: $2x = x + x < x$ \\
\bottomrule
\end{tabular}
\caption{The four cases of one BXH execution.}
\label{tab:cases}
\end{table}

So, whatever the sign, BXH branches exactly when $b = 0$, with one
exception: if $x = 0$, then $2x = x$ and BXH does not branch, although
$b = 0$. BXLE branches in exactly the opposite cases. Because each
execution leaves the shifted value in the register, the next execution tests
the next bit. A sequence of BXH instructions therefore reads the switches
one at a time, from the high-order end.

\section{The switch word}

Two positions in the word are not switches (Figure~\ref{fig:layout}).
\begin{itemize}
\item \textbf{The leading bit} (bit~0) is never tested: the first execution
  tests bit~1. Its value does not matter.
\item \textbf{The trailing bit} (bit~31) must be \textbf{one}. It guarantees
  that the register is never zero while switches remain to be tested, which
  rules out the exception above. With a zero there, a switch that is zero
  and is followed only by zero switches would read as one.
\end{itemize}
That leaves bits 1 to 30 for thirty switches in a fullword. Per the forwarded
note~\cite{metz2024}, the 64-bit forms (BXHG, BXLEG, JXHG and JXLEG) give 62
in a doubleword; Section~7 shows both widths.

\begin{figure}[ht]
\centering
\begin{tikzpicture}[x=7mm,y=7mm]
  \cell{0}{0}{0}{gray!25}
  \foreach \i in {1,...,6} { \cell{\i}{0}{$f_{\i}$}{yellow!25} }
  \node at (7.5,0) {$\cdots$};
  \cell{9}{0}{$f_{29}$}{yellow!25}
  \cell{10}{0}{$f_{30}$}{yellow!25}
  \cell{11}{0}{1}{green!25}
  \node[font=\scriptsize] at (0,0.75) {0};
  \node[font=\scriptsize] at (1,0.75) {1};
  \node[font=\scriptsize] at (10,0.75) {30};
  \node[font=\scriptsize] at (11,0.75) {31};
  \draw (0,-0.6) -- (0,-1.1) node[below, font=\small, align=center]
    {leading bit\\(not tested)};
  \draw (5.5,-0.6) -- (5.5,-1.1) node[below, font=\small] {30 switches};
  \draw (11,-0.6) -- (11,-1.1) node[below, font=\small, align=center]
    {trailing guard\\(must be 1)};
\end{tikzpicture}
\caption{The fullword of switches. Bit numbers are those of the Principles of
Operation, with bit~0 the sign.}
\label{fig:layout}
\end{figure}

\section{A worked example}

Table~\ref{tab:walk} follows a register through six tests, shown at 8 bits
so that every bit is visible: a leading 0, the switches 1~0~1~1~0~0, and the
trailing guard~1. The arithmetic is the same as in 32 bits, modulo $2^8$
instead of $2^{32}$.

\begin{table}[ht]
\centering
\begin{tabular}{clrlrcc}
\toprule
test & $x$ (bits) & $x$ & $2x$ (bits) & $2x$ & BXH & switch \\
\midrule
1 & \texttt{01011001} &   89 & \texttt{10110010} &  $-78$ & falls through & 1 \\
2 & \texttt{10110010} &  $-78$ & \texttt{01100100} &  100 & branches & 0 \\
3 & \texttt{01100100} &  100 & \texttt{11001000} &  $-56$ & falls through & 1 \\
4 & \texttt{11001000} &  $-56$ & \texttt{10010000} & $-112$ & falls through & 1 \\
5 & \texttt{10010000} & $-112$ & \texttt{00100000} &   32 & branches & 0 \\
6 & \texttt{00100000} &   32 & \texttt{01000000} &   64 & branches & 0 \\
\bottomrule
\end{tabular}
\caption{Six successive executions of \texttt{BXH r,r,target} (8-bit
illustration). Wherever the sign changes (tests 1, 2, 3 and 5) the addition
overflows, and the overflow is ignored. In tests 5 and 6 only the trailing
guard keeps the register non-zero.}
\label{tab:walk}
\end{table}

The table was generated by program, and the rule was checked exhaustively
for every switch pattern at register widths of 8, 12 and 16 bits: with the
trailing guard set to one, every switch is read correctly; with it set to
zero, a quarter of all patterns are misread.

\section{In code}

Since BXH branches when the switch is zero, it can branch around the code
for a switch that is off; BXLE reverses the sense, branching when the switch
is on. Every test is executed in turn, whatever the earlier results: unlike a
SELECT, which stops at the first successful test, the sequence must run in
order, because each test depends on the shift performed by the one
before~\cite{ibmmain2024}.

On current machines the relative-branch forms need no base register for the
branch target. JXH and JXLE are the High Level Assembler mnemonics for BRANCH
RELATIVE ON INDEX HIGH and LOW OR EQUAL (BRXH, BRXLE), which operate on 32-bit
values; JXHG and JXLEG (BRXHG, BRXLG) are their 64-bit counterparts. The
arithmetic of the trick is the same at either width. The sections that follow
rest on the forwarded note~\cite{metz2024} and on simulation of the
arithmetic at 32 and 64 bits; the z/Architecture Principles of Operation was
not consulted for this paper.

\subsection{32-bit instructions: 30 switches in a fullword}

AMODE(24) and AMODE(31) are the same case: the 32-bit forms test 30 switches
in either mode. The example shows AMODE(31).

\begin{Verbatim}[fontsize=\small]
SWTEST   CSECT
SWTEST   AMODE 31
SWTEST   RMODE ANY
         L     R15,STATUS         odd register: R1 = R3 = R15
         JXH   R15,R15,SKIP1      switch 1 off: skip action 1
*        ... action 1 ...
SKIP1    JXH   R15,R15,SKIP2      switch 2 off: skip action 2
*        ... action 2 ...
SKIP2    DS    0H
*        ...
STATUS   DC    AL.1(0,FLAG1,...,FLAG30,1)    4 bytes
\end{Verbatim}

With JXLE the sense is reversed, so a test can branch to the code for a
switch that is on:

\begin{Verbatim}[fontsize=\small]
         JXLE  R15,R15,DO1        switch 1 on: go to action 1
\end{Verbatim}

\subsection{Grande instructions: 62 switches in a doubleword}

The grande (64-bit) forms JXHG and JXLEG use the whole 64-bit register. It is
the instruction, not the addressing mode, that sets the width, so no
AMODE(64) is needed: a doubleword holds a leading bit, 62 switches and the
trailing guard. The register is loaded with LG.

\begin{Verbatim}[fontsize=\small]
         LG    R15,STATUS64       odd register: R1 = R3 = R15
         JXHG  R15,R15,SKIP1      switch 1 off: skip action 1
*        ... action 1 ...
SKIP1    JXHG  R15,R15,SKIP2      switch 2 off: skip action 2
*        ... action 2 ...
SKIP2    DS    0H
*        ...
         DS    0D
STATUS64 DC    AL.1(0,FLAG1,...,FLAG62,1)    8 bytes
\end{Verbatim}

As with JXLE, JXLEG reverses the sense:

\begin{Verbatim}[fontsize=\small]
         JXLEG R15,R15,DO1        switch 1 on: go to action 1
\end{Verbatim}

The simulations described in the previous section were repeated at 32 and 64
bits for random switch patterns, weighted towards runs of trailing zero
switches: with the trailing guard set to one, 30 and 62 switches respectively
were read without error.

\section*{Acknowledgments}
The technique and its first description are William Collier's. This paper
was drafted with the assistance of Claude (Anthropic), an AI engine, which
also carried out the simulations; the author is responsible for its content.

\begin{thebibliography}{9}
\bibitem{collier1966}
William Collier, IBM Technical Report TR 00.1412-1, IBM, Poughkeepsie, NY,
22 June 1966. Title not known to the author; cited by number, as Collier
cites it.

\bibitem{poo360}
\emph{IBM System/360 Principles of Operation}, Form A22-6821-6, IBM,
January 1967, pp.~66--66.1, ``Branch On Index High'' and ``Branch On Index
Low or Equal.''
\url{https://bitsavers.org/pdf/ibm/360/princOps/A22-6821-6_360PrincOpsJan67.pdf}

\bibitem{metz2024}
Seymour J. Metz, ``Fw: BXLE instruction,'' ASSEMBLER-LIST, 9 June 2024,
forwarding William Collier's note of 2 November 2022.
\url{https://www.mail-archive.com/assembler-list@listserv.uga.edu/msg16311.html}

\bibitem{ibmmain2024}
``BXLE instruction'' thread, IBM-MAIN, 9 June 2024: Jim Thomas,
\url{https://www.mail-archive.com/ibm-main@listserv.ua.edu/msg138208.html};
Seymour J. Metz,
\url{https://www.mail-archive.com/ibm-main@listserv.ua.edu/msg138211.html};
Gabe Goldberg,
\url{https://www.mail-archive.com/ibm-main@listserv.ua.edu/msg138212.html}.
\end{thebibliography}

\end{document}
